
\documentclass[%
pdftex,
oneside,			% Einseitiger Druck.
11pt,				% Schriftgroesse
parskip=half,		% Halbe Zeile Abstand zwischen Absätzen.
%	topmargin = 10pt,	% Abstand Seitenrand (Std:1in) zu Kopfzeile [laut log: unused]
headheight = 12pt,	% Höhe der Kopfzeile
%	headsep = 30pt,	% Abstand zwischen Kopfzeile und Text Body  [laut log: unused]
headsepline,		% Linie nach Kopfzeile.
footsepline,		% Linie vor Fusszeile.
footheight = 16pt,	% Höhe der Fusszeile
abstracton,		% Abstract Überschriften
DIV=calc,		% Satzspiegel berechnen
BCOR=8mm,		% Bindekorrektur links: 8mm
headinclude=false,	% Kopfzeile nicht in den Satzspiegel einbeziehen
footinclude=false,	% Fußzeile nicht in den Satzspiegel einbeziehen
listof=totoc,		% Abbildungs-/ Tabellenverzeichnis im Inhaltsverzeichnis darstellen
toc=bibliography,	% Literaturverzeichnis im Inhaltsverzeichnis darstellen
]{scrreprt}	% Koma-Script report-Klasse, fuer laengere Bachelorarbeiten alternativ auch: scrbook


\def \titel {Entwicklung eines KI-basierten Assistenzsystems mit transparenter Entscheidungsbegründung zur quantitativen Portfoliooptimierung im Wertpapierhandel}
\def \arbeit {Masterarbeit}
\def \autor {Jasper Bremenkamp}
\def \abschluss {Master of Science}
\def \studiengang {Data Science}
\def \datumAbgabe {\today}
\def \modul {Masterarbeit}
\def \matrikelnr {92125193}
\def \tutor {Christian Müller-Kett}

\def \seitenrand {2cm}

% PDF Einstellungen
\usepackage{hyperref}

\hypersetup{
    unicode = true,
    pdftitle={\titel},
    pdfauthor={\autor},
    pdfsubject={\arbeit},
    pdfcreator={pdflatex, LaTeX with KOMA-Script},
    pdfpagemode=UseOutlines, 		% Beim Oeffnen Inhaltsverzeichnis anzeigen
    pdfdisplaydoctitle=true, 		% Dokumenttitel statt Dateiname anzeigen.
    pdflang={de}, 			% Sprache des Dokuments.
    hidelinks
}

\usepackage[left=2cm,right=2cm,top=2cm,bottom=2cm]{geometry} % margins	% Seitenränder und Abstände
\usepackage[activate]{microtype} %Zeilenumbruch und mehr
\usepackage[onehalfspacing]{setspace}
\usepackage{makeidx}
\usepackage[autostyle=true,german=quotes]{csquotes}
\MakeOuterQuote{"} % Für deutsche Anführungszeichen
\usepackage{longtable}
\usepackage{enumitem}	% mehr Optionen bei Aufzählungen
\usepackage{graphicx}
\usepackage{pdfpages}   % zum Einbinden von PDFs
\usepackage{xcolor} 	% für HTML-Notation
\usepackage{float}
\usepackage{array}
\usepackage{calc}		% zum Rechnen (Bildtabelle in Deckblatt)
\usepackage[right]{eurosym}
\usepackage{wrapfig}
\usepackage{pgffor} % für automatische Kapiteldateieinbindung
\usepackage[perpage, hang, multiple, stable]{footmisc} % Fussnoten
\usepackage[printonlyused]{acronym} % falls gewünscht kann die Option footnote eingefügt werden, dann wird die Erklärung nicht inline sondern in einer Fußnote dargestellt
\usepackage{listings}
\usepackage{bookmark}
\usepackage[english, ngerman]{babel}
\usepackage[utf8]{inputenc}
\usepackage{scrhack}
\usepackage{txfonts} % Use a Times-new-roman open-source clone
\usepackage{comment}

\includecomment{comment} % Kommentare ein und Ausblenden

    %%%%%%% Neues Design %%%%%%%%%%%
    \setkomafont{chapter}{\Large}
    \setkomafont{section}{\large}
    \setkomafont{subsection}{\large}

    %\renewcommand*\chapterheadstartvskip{\vspace*{-0cm}} %Kapitel nach oben verschieben

    \addtolength{\footskip}{-0.7cm}% foot larger by 0,7 cm  (Raises the page number)
    \setlength{\parindent}{6pt} % Indent at start of paragraphs  6pt

    \renewcommand\thesection{\arabic{section}}
    \renewcommand\thefigure{\arabic{figure}}

    %%%%%%%%%%%%%%%%%%%%%%%%%%%%%%%

% Notizen. Einsatz mit \todo{Notiz} oder \todo[inline]{Notiz}.
\usepackage[obeyFinal,backgroundcolor=yellow,linecolor=black]{todonotes}
% Alle Notizen ausblenden mit der Option "final" in \documentclass[...] oder durch das auskommentieren folgender Zeile
% \usepackage[disable]{todonotes}

% der neue Zitierstil ist Käse
\global\def\VAlterZitierstil{true}

\global\def\VTrueValue{true}

\ifx\VAlterZitierstil\VTrueValue
    \usepackage[
        backend=bibtex,     % empfohlen. Falls biber Probleme macht: bibtex
        bibwarn=true,
        bibencoding=utf8,   % wenn .bib in utf8, sonst ascii
        sortlocale=de_DE,
        sorting=none,
        style=numeric-comp
    ]{biblatex}

    \addbibresource{bibliographie.bib}
\else
    % New IUBH Cite Style / Old Cite style below
    \usepackage{natbib}
    \setcitestyle{aysep={}} % remove comma as delimiter
\fi


\begin{document}

\pagenumbering{Roman}

\begin{titlepage}
    \pdfbookmark{\titel}{}
    \begin{longtable}{p{7cm} p{12cm}}
        {\raisebox{\ht\strutbox-\totalheight}{\includegraphics[height=3cm]{assets/iu.png}}}
    \end{longtable}
    \enlargethispage{20mm}
    \begin{center}
        \vspace*{12mm}	{\LARGE\textbf \titel }\\
        \vspace*{12mm}	{\large\textbf \arbeit}\\
        \vspace*{12mm}	Prüfungsleistung für den\\
        \vspace*{3mm}		{\textbf \abschluss}\\
        \vspace*{12mm}	des Studiengangs \studiengang\\
    \vspace*{3mm}		an der Internationalen Hochschule\\
        \vspace*{12mm}	von\\
        \vspace*{3mm}		{\large\textbf \autor}\\
        \vspace*{12mm}	\datumAbgabe
    \end{center}
    \vfill
    \begin{spacing}{1.2}
    \begin{tabbing}
        mmmmmmmmmmmmmmmmmmmmmmmmmm             \= \kill
        \textbf{Kursbezeichnung} \> \modul\\
        \textbf{Matrikelnummer} \> \matrikelnr\\
        \textbf{Tutor} \> \tutor\\
    \end{tabbing}
    \end{spacing}
\end{titlepage}

\newpage

% Inhaltsverzeichnis
\begin{spacing}{1.1}
    \begingroup

        % auskommentieren für Seitenzahlen unter Inhaltsverzeichnis
        \renewcommand*{\chapterpagestyle}{empty}
        \pagestyle{empty}

        \setcounter{tocdepth}{3}
        %für die Anzeige von Unterkapiteln im Inhaltsverzeichnis
        %\setcounter{tocdepth}{2}

        \tableofcontents
        \clearpage
    \endgroup
\end{spacing}
\newpage

%Abbildungsverzeichnis
\cleardoublepage
\listoffigures
\clearpage

\lstlistoflistings %List of listings
\clearpage

%Tabellenverzeichnis
\cleardoublepage
\listoftables
\clearpage

\pagenumbering{arabic}

% Allgemeine To-Dos bzgl. Dokument:
% Chapter in einzelne .tex Files auslagern
% ...

%-------------------------
\chapter{Einleitung}
\begin{comment}
\textbf{To-Dos:}
\begin{itemize}
  \item Prägnante Darstellung des Problems und der Forschungslücke
  \item Klare Formulierung der Forschungsfragen als roter Faden
  \item Abgrenzung des Untersuchungsbereichs ( Fokus auf Transparenz)
  \item Kurzer Ausblick auf den praktischen Beitrag der Arbeit

  \item[\textcolor{blue}{\faLightbulbO}] \textbf{Ideen:}
  \begin{itemize}
    \item Einbettung eines Fallbeispiels (z. B. Portfolio eines Musteranlegers)
    \item Interaktives Diagramm zur Visualisierung der Blackbox-Problematik
    \item Vergleichstabelle: Eigenschaften vs. existierende Robo-Advisor oder andere Systeme
    \item Historischer Exkurs zu Transparenzanforderungen in der Finanzbranche
  \end{itemize}
\end{itemize}
\end{comment}

%-------------------------
\chapter{Stand der Forschung}
\begin{comment}
\textbf{To-Dos:}
\begin{itemize}
  \item Systematische Aufarbeitung existierender Portfolio-Optimierungsansätze
  \item Kritischer Vergleich von Erklärbarkeitsmethoden (SHAP/LIME vs. Alternativen)
  \item Identifikation offener Herausforderungen aktueller Systeme
  \item Theoretische Fundierung der Risikobewertungsmetriken

  \item[\textcolor{blue}{\faLightbulbO}] \textbf{Ideen:}
  \begin{itemize}
    \item Meta-Analyse der Performance kommerzieller Tools
    \item Cross-industrieller Vergleich: Finanz-KI vs. Medizin-KI-Erklärbarkeit
    \item Experteninterview-Zitate zur Validierung der Forschungslücke
    \item Network Graph von Zitationsbeziehungen relevanter Papers
  \end{itemize}
\end{itemize}
\end{comment}

%-------------------------
\chapter{Methodik}
\begin{comment}
\textbf{To-Dos:}
\begin{itemize}
  \item Detaillierte Beschreibung des DSR-Ansatzes mit Design-Zyklen
  \item Nachvollziehbare Begründung der Modellauswahl (z. B, DeepSeek R1) und anderen technischen Gegebenheiten (z. B. Flask, Bootstrap etc)
  \item Dokumentation der Evaluationskriterien vor Implementierung
  \item Klare Zuordnung von Methoden zu Forschungsfragen

  \item[\textcolor{blue}{\faLightbulbO}] \textbf{Ideen:}
  \begin{itemize}
    \item Mockup-Prototypen früher Designiterationen zeigen
    \item Sensitivity-Analyse für Modellhyperparameter
    \item Entscheidungsmatrix zur Framework-Auswahl (Ollama vs. Alternatives)
    \item Live-Demo-Video des Entwicklungsprozesses verlinken
  \end{itemize}
\end{itemize}
\end{comment}

%-------------------------
\chapter{Systemarchitektur und Implementierung}
\begin{comment}
\textbf{To-Dos:}
\begin{itemize}
  \item Modulare Darstellung der Systemkomponenten mit Schnittstellen
  \item Reproduzierbare Code-Beispiele für Kernfunktionen
  \item Diskussion technischer Limitationen (z. B. Latenzzeiten)
  \item Dokumentation der genutzten Frameworks/Tools (z. B. Ollama, Flask)

  \item[\textcolor{blue}{\faLightbulbO}] \textbf{Ideen:}
  \begin{itemize}
    \item \textbf{Privacy-by-Design:} On-Device-Federated-Learning-Ansatz skizzieren
    \item \textbf{Plug-in-System:} Modularer Aufbau für Drittanbieter-Erweiterungen
  \end{itemize}

\end{itemize}
\end{comment}

%-------------------------
\chapter{Evaluation}
\begin{comment}
\textbf{To-Dos:}
\begin{itemize}
  \item Transparente Darstellung der Backtesting-Ergebnisse
  \item Statistische Signifikanzprüfung der Performance
  \item Kritische Reflexion methodischer Limitationen

    \begin{itemize}
    \item \textbf{Stress-Tests:} Simulation von Marktcrashs (z. B. COVID-19, 2008)
    \item \textbf{Erklärungs-Robustheit:} Messung, wie konsistent die KI-Begründungen bei leicht veränderten Inputs sind
  \end{itemize}
\end{itemize}
\end{comment}

%-------------------------
\chapter{Diskussion}
\begin{comment}
\textbf{To-Dos:}
\begin{itemize}
  \item Interpretation der Ergebnisse im Kontext der Forschungsfragen
  \item Einordnung in den aktuellen Forschungsstand
  \item Praktische Implikationen für Finanztechnologie
  \item Offene Herausforderungen und ethische Aspekte

  \item[\textcolor{blue}{\faLightbulbO}] \textbf{Ideen:}
  \begin{itemize}
    \item Szenario-Analyse: Krisenresilienz (2008/2020 Simulation)
  \end{itemize}
\end{itemize}
\end{comment}

%-------------------------
\chapter{Fazit und Ausblick}
\begin{comment}
\textbf{To-Dos:}
\begin{itemize}
  \item Prägnante Zusammenfassung der Kernbeiträge
  \item Beantwortung der Forschungsfragen mit Evidenzbezug
  \item Realistische Skizzierung zukünftiger Forschungsrichtungen
  \item Abschließende Einschätzung der Übertragbarkeit

  \item[\textcolor{blue}{\faLightbulbO}] \textbf{Ideen:}
  \begin{itemize}
    \item \textbf{Cross-Asset Vision:} Übertragung auf Krypto oder andere vergleichbare handelbare Assets
  \end{itemize}

\end{itemize}
\end{comment}



\ifx\VAlterZitierstil\VTrueValue
    \printbibliography
\else
    \bibliographystyle{iubh}
    \bibliography{bibliographie.bib}
\fi
\end{document}
